\chapter{Cơ sở lý thuyết và Nền tảng công nghệ}

\section{Tổng quan về Thương mại điện tử}

Thương mại điện tử là hình thức giao dịch thương mại được thực hiện thông qua mạng Internet và các phương tiện điện tử. Theo nghĩa rộng, TMĐT không chỉ bao gồm giao dịch mua bán trực tuyến mà còn bao trùm toàn bộ chuỗi giá trị từ hoạt động marketing và quảng bá sản phẩm, quản lý đơn hàng và kho hàng, xử lý thanh toán, phối hợp giao vận, đến chăm sóc khách hàng sau bán hàng \cite{chaffey2019digital}.

Mô hình B2C (Business-to-Consumer), trong đó doanh nghiệp bán sản phẩm trực tiếp đến người tiêu dùng cuối, là mô hình phổ biến và được áp dụng trong hệ thống TechStore của khóa luận này. So với kênh bán lẻ truyền thống, mô hình B2C trực tuyến mang lại nhiều lợi thế như tiếp cận thị trường rộng hơn không bị giới hạn địa lý, chi phí vận hành thấp hơn so với cửa hàng vật lý, khả năng cung cấp thông tin sản phẩm phong phú và chi tiết hơn, cùng khả năng thu thập và phân tích dữ liệu hành vi người dùng để cải thiện trải nghiệm mua sắm.

Một hệ thống TMĐT B2C hiện đại được tổ chức theo kiến trúc ba tầng. Tầng giao diện người dùng (Presentation Layer) bao gồm toàn bộ giao diện mà người dùng tương tác trực tiếp, được xây dựng với các công nghệ frontend như React, Vue hay Angular. Tầng xử lý nghiệp vụ (Business Logic Layer) chứa toàn bộ logic kinh doanh như quản lý đơn hàng, tính toán giá, áp dụng mã giảm giá và xử lý thanh toán. Tầng dữ liệu (Data Layer) bao gồm cơ sở dữ liệu quan hệ, hệ thống cache và các dịch vụ lưu trữ.

\section{Kỹ thuật RAG (Retrieval-Augmented Generation)}

\subsection{Giới thiệu và động lực}

Các mô hình ngôn ngữ lớn (Large Language Models, LLM) như GPT-4 hay Gemini sở hữu khả năng sinh văn bản tự nhiên vượt trội nhưng có một hạn chế căn bản: kiến thức của chúng bị giới hạn bởi dữ liệu huấn luyện và không thể cập nhật theo thời gian thực. Điều này đặc biệt problematic trong ứng dụng TMĐT, nơi thông tin về giá cả, tồn kho và sản phẩm mới thay đổi liên tục hàng ngày.

RAG (Retrieval-Augmented Generation) được Lewis và cộng sự giới thiệu năm 2020 \cite{lewis2020rag} như một giải pháp kết hợp tốt nhất của hai thế giới: độ chính xác thông tin từ cơ sở dữ liệu thực tế và khả năng diễn đạt tự nhiên của LLM. Ý tưởng cốt lõi của RAG là trước khi sinh câu trả lời, hệ thống sẽ truy xuất (retrieve) các tài liệu hoặc đoạn thông tin liên quan nhất từ nguồn dữ liệu ngoài, sau đó tăng cường (augment) prompt bằng những thông tin này, và cuối cùng LLM sinh (generate) câu trả lời dựa trên context đã được bổ sung.

\subsection{Kiến trúc RAG}

Kiến trúc RAG bao gồm hai giai đoạn chính. Giai đoạn Indexing (lập chỉ mục) được thực hiện offline, trong đó mỗi tài liệu hoặc sản phẩm trong cơ sở dữ liệu được chuyển đổi thành vector embedding thông qua một mô hình embedding, sau đó vector này được lưu vào vector store. Trong hệ thống TechStore, mỗi sản phẩm được biểu diễn bằng một đoạn văn bản phong phú kết hợp tên sản phẩm, thương hiệu, danh mục, mô tả ngắn và thông tin giá, rồi được nhúng thành vector 1024 chiều.

Giai đoạn Query (truy vấn) được thực hiện theo thời gian thực khi người dùng đặt câu hỏi. Câu hỏi được chuyển đổi thành vector embedding bằng cùng mô hình đã dùng khi lập chỉ mục, sau đó thực hiện tìm kiếm tương đồng trong vector store để truy xuất các tài liệu liên quan nhất. Những tài liệu này được tích hợp vào prompt gửi cho LLM, và LLM sinh câu trả lời dựa trên context đã được bổ sung.

\subsection{Vector Embedding}

Vector embedding là quá trình ánh xạ văn bản thành một điểm trong không gian vector nhiều chiều sao cho các văn bản có nghĩa tương đồng được ánh xạ đến các điểm gần nhau \cite{devlin2019bert}. Chất lượng của embedding quyết định trực tiếp đến khả năng tìm kiếm ngữ nghĩa của hệ thống. Với câu hỏi bằng tiếng Việt về sản phẩm công nghệ, hệ thống cần sử dụng mô hình embedding đa ngôn ngữ có khả năng hiểu cả tiếng Việt lẫn tiếng Anh.

Trong hệ thống TechStore, ba mô hình embedding được sử dụng theo chuỗi ưu tiên (fallback chain). Mô hình ưu tiên nhất là Jina Embeddings v3 \cite{jinaai2024}, một mô hình embedding đa ngôn ngữ mạnh mẽ với 1024 chiều và khả năng xử lý tốt tiếng Việt thông qua API của Jina AI. Nếu Jina không khả dụng, hệ thống chuyển sang mô hình multilingual-e5-large-instruct của HuggingFace với tiền tố hướng dẫn đặc biệt để tối ưu hóa cho bài toán tìm kiếm sản phẩm. Cuối cùng, mô hình multilingual-e5-large đóng vai trò là lựa chọn dự phòng cuối cùng. Cả ba mô hình đều tạo ra vector 1024 chiều đảm bảo tính tương thích khi chuyển đổi giữa các provider.

\subsection{Hybrid Search: kết hợp Semantic và Keyword Search}

Một trong những đóng góp kỹ thuật quan trọng của hệ thống là cơ chế Hybrid Search kết hợp hai phương thức tìm kiếm bổ sung cho nhau. Semantic Search (tìm kiếm ngữ nghĩa) sử dụng cosine similarity để đo độ tương đồng giữa vector query và vector sản phẩm:

\begin{equation}
\text{similarity}(A, B) = \frac{A \cdot B}{\|A\| \cdot \|B\|} = \frac{\sum_{i=1}^{n} A_i B_i}{\sqrt{\sum_{i=1}^{n} A_i^2} \cdot \sqrt{\sum_{i=1}^{n} B_i^2}}
\end{equation}

Semantic Search có ưu điểm là hiểu được các câu hỏi đồng nghĩa và khác cách diễn đạt, ví dụ "điện thoại tốt dưới 10 triệu" và "smartphone giá rẻ chất lượng cao" sẽ cho kết quả tương tự. Tuy nhiên, phương pháp này có thể bỏ sót các kết quả khớp exact match quan trọng, đặc biệt với tên sản phẩm cụ thể như "iPhone 17 Pro Max".

Keyword Search (tìm kiếm từ khóa) được cài đặt theo phong cách BM25 \cite{robertson2009probabilistic}, một thuật toán xác suất xếp hạng tài liệu dựa trên tần suất xuất hiện của từ khóa query trong tài liệu. Trong cài đặt của TechStore, mỗi token trong query được tìm kiếm trong tên sản phẩm (trọng số nhân 3) và trong văn bản mô tả (trọng số nhân 1), sau đó score được điều chỉnh theo tỷ lệ token khớp so với tổng số token trong query. Keyword Search đặc biệt hiệu quả với tên sản phẩm, mã model và số thế hệ cụ thể.

Hai phương thức này được kết hợp theo cơ chế fusion: các sản phẩm xuất hiện trong cả kết quả semantic và keyword được tăng điểm thêm 0.05 để ưu tiên. Các sản phẩm chỉ xuất hiện trong kết quả keyword (semantic bỏ sót) được inject vào danh sách cuối với điểm tính toán từ keyword score nhân hệ số điều chỉnh, kèm flag \texttt{lowConfidence} để LLM biết và diễn giải phù hợp. Kết quả cuối được sắp xếp theo score tổng hợp giảm dần và lấy top-N sản phẩm.

\section{Kiến trúc Modular Monolith}

\subsection{Khái niệm và đặc điểm}

Kiến trúc Modular Monolith là hướng tiếp cận thiết kế hệ thống kết hợp ưu điểm của cả Monolith và Microservices trong khi hạn chế nhược điểm của cả hai \cite{newman2019monolith}. Trong kiến trúc này, toàn bộ ứng dụng vẫn được triển khai như một đơn vị duy nhất (monolith) nhưng được tổ chức nội bộ thành các module độc lập rõ ràng (modular) với ranh giới rõ ràng giữa các module.

So với Microservices, Modular Monolith đơn giản hơn đáng kể trong triển khai (không cần container orchestration, service mesh hay distributed tracing), loại bỏ overhead của network calls giữa các service và phù hợp hơn với quy mô nhóm nhỏ và giai đoạn đầu phát triển sản phẩm. So với Monolith truyền thống, Modular Monolith duy trì được tính tách biệt về trách nhiệm, dễ bảo trì và về lý thuyết có thể tách các module thành microservices độc lập khi cần mở rộng quy mô trong tương lai.

\subsection{Vertical Slice và Dependency Injection}

Mỗi module trong hệ thống TechStore được tổ chức theo pattern Vertical Slice, nghĩa là mỗi module đóng gói đầy đủ tất cả các tầng cần thiết cho một domain nghiệp vụ cụ thể, từ HTTP routing đến business logic và data access. Cụ thể, mỗi module bao gồm file \texttt{module.js} đóng vai trò DI factory nhận dependencies qua constructor và khởi tạo toàn bộ object graph nội bộ, file \texttt{routes.js} định nghĩa các HTTP endpoints và middleware chain, thư mục \texttt{controllers/} xử lý HTTP layer, thư mục \texttt{services/} chứa business logic, thư mục \texttt{repositories/} thực hiện data access qua Sequelize và thư mục \texttt{validators/} định nghĩa Zod schemas.

Dependency Injection (DI) là kỹ thuật thiết kế trong đó các phụ thuộc của một đối tượng được cung cấp từ bên ngoài thay vì đối tượng tự tạo ra. Trong hệ thống TechStore, file \texttt{src/app.js} đóng vai trò DI container duy nhất, khởi tạo tất cả shared resources (Sequelize models, EventBus, emailService) và inject chúng vào các module factory function. Điều này đảm bảo mỗi module không tự import models hay services trực tiếp, giúp dễ dàng thay thế implementation (đặc biệt quan trọng khi unit test với mock dependencies).

\subsection{EventBus và Event-Driven Communication}

EventBus là thành phần quan trọng cho phép các module giao tiếp bất đồng bộ mà không tạo ra sự phụ thuộc trực tiếp giữa chúng. Trong hệ thống TechStore, EventBus được cài đặt như một singleton in-process pub/sub, nghĩa là tất cả modules chia sẻ cùng một instance EventBus trong một process Node.js duy nhất. Một module có thể publish event với type và payload, trong khi các module khác đã subscribe handler cho event type đó sẽ tự động được gọi. Cơ chế \texttt{Promise.allSettled} đảm bảo một handler bị lỗi không ảnh hưởng đến các handler khác.

Ví dụ điển hình là luồng tạo đơn hàng: module \texttt{orders} publish event \texttt{order.created} sau khi tạo đơn thành công, module \texttt{inventory} subscribe handler cho event này để ghi audit log thay đổi tồn kho. Tương tự, khi đơn hàng bị hủy, event \texttt{order.cancelled} được publish và inventory ghi log restore hàng về kho. Cách thiết kế này giúp module \texttt{orders} không cần biết đến sự tồn tại của module \texttt{inventory}.

\section{Công nghệ Backend}

\subsection{Node.js và Express.js}

Node.js là nền tảng thực thi JavaScript phía máy chủ, được xây dựng trên V8 JavaScript Engine của Google. Điểm mạnh then chốt của Node.js là mô hình event-driven, non-blocking I/O, cho phép xử lý hàng nghìn kết nối đồng thời mà không tạo thread riêng cho mỗi kết nối. Đây là lợi thế quan trọng với API server TMĐT, nơi phần lớn thời gian là chờ đợi kết quả từ database, cache hay external services.

Hệ thống TechStore sử dụng Node.js phiên bản 20 LTS, phiên bản ổn định nhất tại thời điểm phát triển, với các tính năng hiện đại như native \texttt{fetch}, \texttt{AbortController} cải tiến và Web Crypto API. Express.js phiên bản 4.18 được dùng làm web framework chính nhờ tính tối giản và hệ sinh thái middleware phong phú. Stack middleware trong \texttt{app.js} được tổ chức cẩn thận: Helmet bảo vệ HTTP headers, CORS kiểm soát cross-origin requests, Morgan ghi log requests, rate limiters bảo vệ chống abuse, sanitizeHtml lọc XSS, compression giảm kích thước response.

\subsection{Sequelize ORM và MySQL 8}

Sequelize là ORM (Object-Relational Mapping) cho Node.js hỗ trợ MySQL, PostgreSQL và SQLite. ORM cho phép làm việc với cơ sở dữ liệu quan hệ thông qua các đối tượng JavaScript thay vì viết SQL thuần, tự động tạo câu truy vấn an toàn tránh SQL Injection, cung cấp system migration để quản lý thay đổi schema theo thời gian và hỗ trợ đầy đủ các tính năng của MySQL như transactions, foreign key constraints và soft delete.

MySQL 8 được chọn làm hệ quản trị cơ sở dữ liệu vì hỗ trợ đầy đủ InnoDB storage engine với ACID transactions, full-text search, JSON column type và SELECT FOR UPDATE cho row-level locking. Hệ thống sử dụng charset utf8mb4\_unicode\_ci để đảm bảo hỗ trợ đầy đủ Unicode bao gồm tiếng Việt với các dấu thanh đầy đủ và emoji. Timezone được cấu hình là +07:00 (múi giờ Việt Nam) để đảm bảo consistency giữa server và database.

\subsection{Redis và Chiến lược Caching}

Redis là hệ thống lưu trữ key-value in-memory tốc độ cao, được sử dụng trong TechStore cho ba mục đích. Mục đích quan trọng nhất là JWT blacklist: khi người dùng đăng xuất, access token của họ được lưu vào Redis với key là JWT ID (jti) và TTL bằng thời gian sống còn lại của token. Mọi request tiếp theo đến middleware authenticate sẽ kiểm tra xem jti có trong blacklist không trước khi chấp nhận. Mục đích thứ hai là caching response của chatbot: kết quả truy vấn chatbot cho intent \texttt{product\_search} được cache trong Redis với TTL 5 phút sử dụng shared cache key (không phụ thuộc userId) để tối đa hóa cache hit rate. Mục đích thứ ba là caching danh sách sản phẩm, danh mục và thương hiệu cho catalog API với TTL 10 đến 30 phút.

Một thiết kế quan trọng là hệ thống không bao giờ crash khi Redis không khả dụng. Config Redis được xây dựng với cơ chế auto-fallback: nếu kết nối Redis thất bại (timeout 1500ms, không retry), hệ thống tự động chuyển sang sử dụng in-memory Map với interface tương đương bao gồm đầy đủ các methods \texttt{setEx}, \texttt{get}, \texttt{del}, \texttt{keys}. Đây là quyết định thiết kế quan trọng đảm bảo tính khả dụng của hệ thống trong mọi tình huống.

\subsection{Xác thực JWT và Google OAuth}

Xác thực người dùng trong TechStore sử dụng hai cơ chế bổ sung cho nhau. JWT (JSON Web Token) \cite{jones2015jwt} là chuẩn mở RFC 7519 định nghĩa cách truyền tải thông tin xác thực dưới dạng JSON object được ký số. Hệ thống sử dụng chiến lược dual-token: access token với TTL 15 phút được lưu trong bộ nhớ frontend (sessionStorage) để bảo mật, và refresh token với TTL 30 ngày được lưu trong httpOnly cookie để tránh JavaScript access. Khi access token hết hạn, frontend tự động gọi endpoint refresh để lấy cặp token mới.

Đặc biệt, hệ thống cài đặt cơ chế token family rotation để phát hiện và ngăn chặn token reuse attack: mỗi refresh token thuộc về một family ID, khi rotate thành công token cũ bị revoke. Nếu một refresh token đã bị revoke được sử dụng lại (dấu hiệu token bị đánh cắp), hệ thống tự động revoke toàn bộ family, buộc user đăng nhập lại trên tất cả thiết bị.

Google OAuth 2.0 được tích hợp cho phép đăng nhập bằng tài khoản Google. Hệ thống hỗ trợ hai chế độ xác thực: xác thực \texttt{id\_token} (mobile flow sử dụng Google Sign-In SDK) và xác thực \texttt{access\_token} (web flow sử dụng Google OAuth popup). Nếu email Google đã tồn tại trong hệ thống, tài khoản sẽ được tự động liên kết; nếu chưa tồn tại, tài khoản mới được tạo tự động.

\subsection{Bảo mật và Rate Limiting}

Bảo mật là yêu cầu thiết yếu của mọi hệ thống TMĐT. TechStore cài đặt nhiều lớp bảo mật. Về phía bảo vệ XSS, thư viện \texttt{sanitize-html} được dùng như middleware toàn cục để loại bỏ tất cả HTML tags và attributes khỏi request body trước khi data đến các handlers. Về phía bảo vệ CSRF, hệ thống kiểm tra Origin header cho tất cả state-changing requests (POST, PUT, PATCH, DELETE) và từ chối những request có Origin không nằm trong whitelist. Content Security Policy được cấu hình qua Helmet để kiểm soát nguồn tải script và style.

Rate limiting được triển khai với 5 limiters có ngưỡng khác nhau: \texttt{apiLimiter} (100 request mỗi 15 phút cho toàn bộ API trong production), \texttt{authLimiter} (10 request mỗi 60 phút cho auth endpoints để chống brute force), \texttt{otpLimiter} (5 request mỗi 15 phút cho OTP endpoints), \texttt{chatbotLimiter} (20 request mỗi 60 giây cho chatbot), và \texttt{chatLimiter} (30 request mỗi 5 phút cho lịch sử chat). Tất cả limiters đều sử dụng \texttt{ProxyStore} tự động upgrade từ in-memory lên Redis khi kết nối thành công.

\section{Công nghệ Frontend}

\subsection{React 18 và TypeScript}

React 18 là phiên bản mới nhất của thư viện xây dựng giao diện người dùng của Meta, mang lại nhiều cải tiến quan trọng so với phiên bản trước. Concurrent Rendering cho phép React ưu tiên render các cập nhật quan trọng hơn, Suspense hỗ trợ lazy loading component và data fetching khai báo tường minh, trong khi automatic batching gom nhóm nhiều state updates vào một lần re-render duy nhất để tối ưu hiệu năng.

TypeScript bổ sung hệ thống kiểu tĩnh vào JavaScript, giúp phát hiện lỗi tại compile time thay vì runtime, cải thiện khả năng bảo trì code thông qua type inference và IDE support, và đặc biệt quan trọng với dự án quy mô lớn như TechStore với 14 features và hàng trăm components.

\subsection{Vite và Feature-Based Architecture}

Vite 5 là build tool thế hệ mới sử dụng native ES Modules trong development mode, cung cấp HMR (Hot Module Replacement) gần như tức thì nhờ chỉ transform các modules thực sự thay đổi. So với Webpack, Vite khởi động nhanh hơn đáng kể (dưới 1 giây so với nhiều giây với Webpack), và production build cũng nhanh hơn nhờ sử dụng Rollup được tối ưu.

Frontend TechStore áp dụng Feature-Based Architecture: mỗi feature trong \texttt{src/features/<name>/} là một unit cô lập hoàn toàn, chứa đầy đủ logic API, components, pages, hooks và types riêng. Không có cross-feature imports trực tiếp (ngoại trừ một trường hợp được phép là \texttt{orders} import \texttt{ReviewModal} từ \texttt{reviews}). Code dùng chung đặt tại \texttt{src/components/}, \texttt{src/stores/}, \texttt{src/hooks/} và \texttt{src/utils/}.

\subsection{TanStack Query v5 và Zustand v5}

TanStack Query (trước đây là React Query) v5 quản lý toàn bộ server state, tức là dữ liệu đến từ API. Thư viện này cung cấp cơ chế caching thông minh với staleTime và gcTime cấu hình được, tự động refetch khi data trở nên stale hoặc window được focus lại, và hỗ trợ optimistic updates cho trải nghiệm người dùng mượt mà hơn. Query key được tổ chức theo hierarchy giúp invalidation chính xác: thay đổi một sản phẩm chỉ invalidate cache liên quan đến sản phẩm đó, không ảnh hưởng toàn bộ cache.

Zustand v5 kết hợp Immer middleware quản lý client state, tức là trạng thái UI không lưu trên server. Hệ thống có 6 Zustand stores: \texttt{auth-store} cho thông tin xác thực và user, \texttt{cart-store} cho giỏ hàng local (persist qua localStorage), \texttt{chat-store} cho lịch sử hội thoại chatbot, \texttt{catalog-store} cho danh sách sản phẩm đã xem gần đây và compare list, \texttt{ui-store} cho theme và notifications, và \texttt{wishlist-store} cho danh sách yêu thích. Mỗi store chỉ lưu những gì thực sự cần thiết ở client side, tránh duplicate với server state được quản lý bởi TanStack Query.

\subsection{Ant Design và Tailwind CSS}

Ant Design v5 là thư viện UI component phong phú của Alibaba, được sử dụng chủ yếu cho phần quản trị admin với các components phức tạp như Table với sorting và filtering, Form với validation tích hợp, Modal, Pagination và DatePicker. Tailwind CSS v3 là utility-first CSS framework được sử dụng cho phần giao diện người dùng public, kết hợp với SCSS cho các global styles, CSS variables cho design tokens và Framer Motion cho animations.

\section{Hệ thống Thanh toán}

\subsection{Kiến trúc tích hợp cổng thanh toán}

Cả MoMo và VNPay đều sử dụng cùng một mô hình tích hợp: frontend gọi backend để tạo payment URL với chữ ký HMAC, backend trả URL về, frontend redirect người dùng đến trang thanh toán của cổng. Sau khi người dùng hoàn thành thanh toán, cổng gửi IPN (Instant Payment Notification) là một HTTP request đến backend server để thông báo kết quả, đồng thời redirect người dùng về frontend qua return URL. Điểm quan trọng là state change (cập nhật trạng thái đơn hàng) PHẢI xảy ra trong IPN handler phía server, không được dựa vào return URL phía client vì client-side redirect có thể bị bypass.

\subsection{MoMo Payment Integration}

MoMo là ví điện tử phổ biến nhất tại Việt Nam với hơn 31 triệu người dùng. Tích hợp MoMo sử dụng HMAC-SHA256 để ký request và xác thực IPN callback. Chuỗi ký bao gồm toàn bộ các trường quan trọng nối theo thứ tự quy định của MoMo như \texttt{accessKey}, \texttt{amount}, \texttt{orderId} và \texttt{requestType}. Để đảm bảo tính duy nhất của mỗi lần thử thanh toán (MoMo yêu cầu orderId unique), hệ thống append timestamp vào orderId gốc: \texttt{ORD-2605-xxx-{timestamp6}}. IPN từ MoMo được gửi qua POST request, và xác thực chữ ký sử dụng \texttt{crypto.timingSafeEqual} thay vì so sánh string thông thường để chống timing attack.

\subsection{VNPay Payment Integration}

VNPay là cổng thanh toán hỗ trợ thẻ ngân hàng nội địa và quốc tế, sử dụng HMAC-SHA512 với quy tắc encoding đặc biệt. Amount được nhân 100 vì VNPay làm việc ở đơn vị xu. Một điểm khác biệt quan trọng với MoMo là VNPay IPN là HTTP GET request (không phải POST), và trong một số trường hợp VNPay yêu cầu return URL cũng phải xử lý state change (cập nhật đơn hàng). Hệ thống xử lý cả hai: IPN handler cho server-to-server notification và return URL handler cho trường hợp IPN chậm hoặc không đến.

\section{Các Công nghệ Hỗ trợ}

\subsection{Nodemailer và Gmail SMTP}

Hệ thống gửi email transactional qua Nodemailer với Gmail SMTP (port 587, STARTTLS) cho sáu loại email: OTP xác thực tài khoản, OTP reset mật khẩu, xác nhận đặt hàng thành công, cập nhật trạng thái đơn hàng, thông báo hủy đơn và thông báo feedback đến admin. Email được gửi bất đồng bộ (fire-and-forget) với Nodemailer pool 3 connections để tránh overhead tạo kết nối mới mỗi lần gửi. Nội dung email sử dụng HTML template với input escaping để tránh XSS trong nội dung email.

\subsection{Sharp và Multer}

Sharp là thư viện xử lý ảnh cho Node.js với hiệu năng cao. Trong TechStore, Sharp được dùng để resize ảnh sản phẩm về kích thước tối đa 800x800 pixels (giữ tỷ lệ), chuyển đổi sang định dạng WebP với quality 85 để giảm kích thước file, strip toàn bộ EXIF metadata (bao gồm thông tin GPS location của ảnh chụp từ điện thoại) và tạo thumbnails với ba kích thước 150x150, 300x300 và 600x600 pixels. Multer xử lý multipart/form-data upload với validation MIME type và giới hạn kích thước file 5MB.

\subsection{Zod và Winston}

Zod là thư viện schema validation TypeScript-first, được dùng để validate toàn bộ input từ client trước khi đến business logic. Điểm mạnh của Zod là tự động strip các field không khai báo trong schema (bảo vệ chống mass assignment), cung cấp type inference tự động cho TypeScript và cho phép compose schema phức tạp từ các primitive đơn giản. Winston là logging framework cho Node.js với hỗ trợ nhiều transport và định dạng log: console với format màu trong development, file JSON trong production với rotation 10MB mỗi file.

\section*{Kết luận chương 2}
\addcontentsline{toc}{section}{Kết luận chương 2}

Chương 2 đã trình bày toàn bộ nền tảng lý thuyết và công nghệ của hệ thống TechStore. Kỹ thuật RAG với Hybrid Search là nền tảng cho tính năng cốt lõi của hệ thống, kiến trúc Modular Monolith đảm bảo cấu trúc phần mềm rõ ràng và dễ bảo trì, trong khi stack công nghệ Node.js + React + MySQL 8 được lựa chọn dựa trên sự trưởng thành, hệ sinh thái phong phú và phù hợp với bài toán TMĐT hiện đại. Các quyết định kiến trúc như auto-fallback Redis, token family rotation và Hybrid Search đều xuất phát từ yêu cầu thực tiễn về tính khả dụng, bảo mật và chất lượng kết quả tìm kiếm.
